\subsubsection{Manifeste}

Die Objekte von Kubernetes, wie beispielsweise Pods, Deployments und ConfigMaps, werden mithilfe von Manifesten festgelegt. 
Diese Manifeste sind normalerweise YAML-Dateien, die den Aufbau, das Verhalten, die Befugnisse und auch den Inhalt der Objekte beschreiben. 
Kubernetes selber sagt das für Manifeste auch JSON genommen werden kann \parencite{kubernetesObjects}, dies ist jedoch eher unüblich.
Manifeste unterscheiden sich je nach Objekt, dennoch gibt es grundlegende Informationen, die angegeben werden müssen.
Dazu gehört, welches Objekt mit Kind erstellt werden soll und welche Version des Objekts mit Apiversion genutzt werden soll. 
Zusätzlich müssen grundlegende Informationen über das Objekt unter Metadaten beschrieben werden, wobei mindestens der Name des Objekts angegeben werden muss. 
Danach kann unter „spec” der gewünschte Zustand und das Verhalten des Objekts festgelegt werden.

\begin{lstlisting}[language=yaml, caption={Beispiel von einen Manifest}]
apiVersion: v1
kind: Pod
metadata:
    name: my-pod
spec:
    ...
\end{lstlisting}